\section{$Z$-scan: hydrogenic-ion $g$-factor across $Z = 2, 6, 8, 14, 50$ (PR \#87)}
\label{sec:eq-zscan-fit}

Source (consumed from branch \texttt{82-hydrogenic-z-scan-g-factor}):
\texttt{Mathematica\_Notebooks/Quantum\_Mechanics/r\_e\_Zscan\_fit.wl}
(157 lines).
Companion to
\texttt{Bethe\_Salpeter/14\_HydrogenicIon\_Zscan.md}.

\paragraph{Setup.}
For each ion of charge $Z$ with measured bound-electron $g$-factor
$g_{\rm meas}(Z)$, test two readings of (III.22)
$g_r(x) = 2\bigl(1 - 4/(2x + 1)\bigr)$:
\begin{itemize}[topsep=2pt,itemsep=0pt]
  \item \textbf{(Z-i) universal cutoff:} $x$ fixed at $x_{\rm univ} =
        0.4994205099128317$ (PR~\#62 triangulated). Framework predicts the
        \emph{same} $g = -2.00231930\dots$ at every $Z$.
  \item \textbf{(Z-ii) Z-scaled cutoff:} invert $g_r$ at each measured $g$
        to get $x^{(Z)} = (2 - a)/(2(2 + a))$ with $a = (|g_{\rm meas}| - 2)/2$;
        fit $x^{(Z)}$ as $c_0 + c_2 (Z\alpha)^2 + c_4 (Z\alpha)^4$.
\end{itemize}

\paragraph{Measurement provenance.}
\begin{center}
\begin{tabular}{lllll}
\toprule
$Z$ & ion             & measured $|g|$ & source                  & DOI \\
\midrule
2  & $^3$He$^+$       & 2.0021774158  & Schneider 2022 \emph{Nature} 606, 878 & 10.1038/s41586-022-04761-7 \\
6  & $^{12}$C$^{5+}$  & 2.0010415902  & Sturm 2014 \emph{Nature} 506, 467     & 10.1038/nature13026 \\
8  & $^{16}$O$^{7+}$  & 2.0000470254  & Verd\'u 2004 \emph{PRL} 92, 093002    & 10.1103/PhysRevLett.92.093002 \\
14 & $^{28}$Si$^{13+}$& 1.9953489587  & Sturm 2011 \emph{PRL} 107, 023002     & 10.1103/PhysRevLett.107.023002 \\
50 & $^{118}$Sn$^{49+}$& 1.9105620590 & Morgner 2023 \emph{Nature} 622, 53    & 10.1038/s41586-023-06453-2 \\
\bottomrule
\end{tabular}
\end{center}

\paragraph{S2 --- per-$Z$ back-fit cutoff $x^{(Z)} = r_e^{(Z)}/r_0$.}
\wolframcheck{Wolfram MCP, executed 2026-05-27, STATE.md iter 11.}
\begin{center}
\begin{tabular}{rrrr}
\toprule
$Z$ & $a_{\rm bound}$       & $x^{(Z)}$           & shift from $x_{\rm univ}$ \\
\midrule
2  & $+0.00108870790$      & $0.49945594221$     & $+3.5\times 10^{-5}$ \\
6  & $+0.00052079510$      & $0.49973967024$     & $+3.2\times 10^{-4}$ \\
8  & $+0.00002351270$      & $0.49998824379$     & $+5.7\times 10^{-4}$ \\
14 & $-0.00232552065$      & $0.50116411391$     & $+1.7\times 10^{-3}$ \\
50 & $-0.04471897050$      & $0.52287086604$     & $+2.3\times 10^{-2}$ \\
\bottomrule
\end{tabular}
\end{center}

\paragraph{S3 --- (Z-i) universal-cutoff test.}
With $x = x_{\rm univ}$ fixed (0 free parameters), the framework predicts
$g_{\rm pred} = 2.00231930436256$ at every $Z$. Residuals:
\[
\{-1.42\times 10^{-4},\; -1.28\times 10^{-3},\; -2.27\times 10^{-3},\;
  -6.97\times 10^{-3},\; -9.18\times 10^{-2}\}
\]
in $\sigma$-units of measurement precision:
\[
\{-3.15\times 10^{5},\; -4.26\times 10^{7},\; -4.94\times 10^{5},\;
  -6.97\times 10^{6},\; -9.18\times 10^{7}\}.
\]
\[
\boxed{\;
  \chi^2_{\rm (Z\text{-}i)} = 1.0282\times 10^{16}
  \;\Longrightarrow\; \text{OUTCOME A DECISIVELY REJECTED.}
\;}
\]

\paragraph{S4 --- (Z-ii) Z-scaling fit.}
Fit of $x^{(Z)}$ vs.~$(Z\alpha)^2$:
\begin{align*}
\text{linear:}\quad x &= 0.499\,383\,5903 + 0.176\,393\,13\,(Z\alpha)^2 \\
\text{quadratic:}\quad x &= 0.499\,420\,6084 + 0.166\,275\,68\,(Z\alpha)^2
                          + 0.074\,154\,07\,(Z\alpha)^4
\end{align*}
\textbf{The quadratic intercept} $c_0 = 0.4994206084$ matches
$x_{\rm univ} = 0.4994205099$ to $9.8\times 10^{-8}$.
\textbf{The quadratic slope} $c_2 = 0.166276$ matches QED's
$1/6 = 0.166667$ to $2.3\times 10^{-3}$.

\paragraph{S5 --- QED bound-state cross-check.}
$|g(Z\alpha)|^{\rm QED} \approx 2\bigl[1 - \tfrac{1}{3}(Z\alpha)^2
- \tfrac{1}{12}(Z\alpha)^4\bigr] + \alpha/\pi$.
The $x^{(Z)}$ back-fit curve mirrors this; the intercept-shift from $1/2$ is
the free anomaly.

\paragraph{Outcome C (verbatim from notebook).}
\begin{itemize}[topsep=2pt,itemsep=0pt]
  \item Outcome~A (single universal cutoff) is excluded at
        $\chi^2 \sim 10^{16}$ over 5 ions ($Z = 2 \ldots 50$).
  \item The (Z-ii) back-fit $x^{(Z)}$ is well-described by
        $x = c_0 + c_2 (Z\alpha)^2 + c_4 (Z\alpha)^4$ with $c_0 =
        x_{\rm univ}$ recovered as $Z\to 0$ to $10^{-7}$, and $c_2 \approx 1/6$
        --- the QED bound-state leading coefficient.
  \item The clean $(Z\alpha)^2$ form is a property of QED's bound-state
        $g(Z\alpha) = 2[1 - (1/3)(Z\alpha)^2 - \dots]$, \emph{inherited} by the
        inversion --- not derived from the framework apparatus (which supplies
        only $g_r(x)$ and leaves each state's cutoff free; the $-1/3$
        coefficient is QED's, not the dual framework's). Therefore
        \textbf{Outcome~C}, extending PR~\#70's lepton-axis verdict
        (``particle-specific through $a_\ell$'') to the $Z$-axis
        (``$Z$-specific through $a_e^{\rm bound}(Z\alpha)$'').
\end{itemize}

\authorreview{The PR~\#87 verdict claims the cutoff is per-$Z$-specific via
QED-inheritance (Outcome~C), and that the (Z-i) universal-cutoff hypothesis
is ruled out at $\chi^2 \sim 10^{16}$. PRs \#84/\#85/\#86 claim ``Branch~A by
construction'' from the framework's algebra (which the present notebook does
not dispute; the algebra has no $Z$-dependence). The resolution PR~\#87
offers: the algebra cannot produce the binding correction, so the framework
\emph{inherits} the QED coefficient via the back-fit, but does not
\emph{derive} it. Does this reading --- Outcome~C as ``QED-inheritance,
framework does not internally derive the binding coefficient'' --- match
your reading of (III.23) leaving the cutoff free per state? Is the muon
cross-check (PR~\#70) the analogous lepton-axis verdict the framework
inherits?}
